% Fiche Verilog / SystemVerilog — recto-verso A4 paysage
% Source : VERILOG-CHEATSHEET-.pdf + AED_SV_UVM_Cheatsheet p. 2-3 (repris, traduit, corrigé)
% Forme une paire avec ../UVM/uvm_cheat.tex : ici le langage, la conception et la
%   simulation ; la vérification (classes, aléatoire, couverture, assertions, UVM) est là-bas.
% Compiler : make (ici ou à la racine), ou xelatex verilog_cheat.tex — une seule passe suffit
% Apparence : ../librairie (fiche.sty + theme-verilog.tex). Macros de contenu :
%   colonne[titre][couleur] (3 par page, séparées par \colbreak) · bloc{titre}{catégorie}
%   catégories : syn typ ctl hw tb agg (couleurs et icônes dans le thème)
%   code · \V|...| code en ligne · «...» dans le code = à remplir
%   \sub sous-titre · \note remarque · \piege piège · \sv badge SystemVerilog · schema
%   optab étiquette | texte · tab3 étiquette | code | remarque · grille{n} · typetab
%   (largeurs mesurées par la librairie : celles écrites entre [] et {} sont ignorées)
% Le code est en police à chasse fixe : les commentaires sont alignés à la main
% (attention : les « » ne sont pas affichés, ils ne comptent pas dans l'alignement).
% Après compilation : grep "\[fiche\]" verilog_cheat.log (place libre, collisions).
\documentclass{article}
\usepackage[verilog]{../librairie/fiche}   % apparence : ../librairie (fiche.sty + theme-verilog.tex)
% --- Langage SystemVerilog pour listings ---
% Classes de mots-clés : [1] rose · [2] cyan · [3] vert · [4] couleur du thème
% (styles dans ../librairie/fiche.sty). Chargé par verilog_cheat.tex.
\lstdefinelanguage{CODE}{
  sensitive=true, alsoletter={\$`},
  morekeywords=[1]{module,endmodule,input,output,inout,parameter,localparam,assign,
    always,always_ff,always_comb,always_latch,initial,begin,end,if,else,case,casez,
    casex,endcase,default,for,foreach,repeat,while,do,forever,fork,join,join_none,
    join_any,task,endtask,interface,endinterface,modport,clocking,endclocking,
    mailbox,semaphore,process,new,function,endfunction,return,generate,endgenerate,genvar,
    posedge,negedge,or,wait,disable,virtual,inside,typedef,enum,struct,union,packed,class,automatic,void},
  morekeywords=[2]{wire,reg,logic,bit,byte,shortint,int,longint,integer,real,
    shortreal,realtime,time,string,event,supply0,supply1,tri,triand,trior,tri0,
    tri1,wand,wor,signed,unsigned},
  morekeywords=[3]{\$display,\$monitor,\$finish,\$stop,\$time,\$random,\$urandom,
    \$readmemh,\$readmemb,\$signed,\$unsigned,
    `timescale,`define,`include},
  morecomment=[l]{//}, morecomment=[s]{/*}{*/}, morestring=[b]",
}
   % coloration syntaxique : \lstdefinelanguage{CODE}

\begin{document}
%%%%%%%%%%%%%%%%%%%%%%%%%%%%%%%%%%%%%%%%%%%%%%%%%%%%%%%%%%%%%%%%%%%%%%%%%%%%%%
%                         RECTO : LE LANGAGE
%%%%%%%%%%%%%%%%%%%%%%%%%%%%%%%%%%%%%%%%%%%%%%%%%%%%%%%%%%%%%%%%%%%%%%%%%%%%%%
\entete{Verilog / SystemVerilog}{recto · le langage}
\noindent
% ======================= colonne 1 : syntaxe, types Verilog =======================
\begin{colonne}[opérateurs · types][col-a]
\begin{bloc}{Opérateurs}{syn}
\begin{tab3}[17mm]{22mm}
Arithm.     & \V|+ - * / %|   & \V|**| puissance \\
Relation.   & \V|> < >= <=|   & résultat sur 1 bit \\
Égalité     & \V|== !=|       & x ou z $\Rightarrow$ résultat x \\
            & \V|=== !==|     & compare aussi x et z \\
            & \V|==? !=?| \sv & x / z à droite = joker \\
Logiques    & \V/&& || !/     & résultat sur 1 bit \\
Bit à bit   & \V/& | ^ ~ ~^/  & et, ou, xor, non, xnor \\
Réduction   & \V/&a |a ^a/    & tous les bits $\to$ 1 bit \\
            & \V/~&a ~|a ~^a/ & \V|^a| = parité \\
Décalages   & \V|<< >>|       & logiques (insère des 0) \\
            & \V|<<< >>>|     & arithm. (\V|>>>| garde le signe) \\
Condition   & \V|c ? a : b|   & a si c vrai, sinon b \\
Ensemble    & \V|a inside {1,4}| & \sv\ appartient à l'ensemble \\
Incrément   & \V|i++ i-- +=|  & \sv\ comme en C \\
Concat.     & \V|{a, b[2:0]}| & assemble les bits \\
Réplication & \V|{4{2'b01}}|  & = \V|8'b01010101| \\
Affectation & \V|=|  \V|<=|   & bloquante · non bloquante \\
Sélection   & \V|v[7:4]|      & intervalle de bits \\
            & \V|v[i+:4]|     & 4 bits à partir de i
\end{tab3}
\end{bloc}

\begin{bloc}{Littéraux}{syn}
\begin{tab3}[17mm]{22mm}
Format      & \V|8'hFF|        & \V|<taille>'<base><valeur>| \\
Bases       & \V|'b 'o 'd 'h|  & binaire, octal, décimal, hexa \\
Signé       & \V|4'sb1010|     & \V|s| = signé (vaut $-6$) \\
Défaut      & \V|1973  -123|   & décimal, 32 bits, signé \\
Sans taille & \V|'hcafe|       & 32 bits \\
Séparateur  & \V|16'hBE_EF|    & \V|_| ignoré (lisibilité) \\
x / z       & \V|4'b10xz|      & inconnu / haute impédance \\
Remplissage & \V|'0 '1 'x 'z|  & \sv\ tous les bits à la valeur
\end{tab3}
\end{bloc}

\begin{bloc}{Types Verilog}{typ}
\textbf{Logique 4 valeurs :} \V|0| \V|1| \V|x| inconnu · \V|z| haute impédance
\begin{tab3}[17mm]{22mm}
Nets        & \V|wire tri|       & fil · bus trois états \\
            & \V|wand wor|       & ET / OU câblé \\
            & \V|tri0 tri1|      & tiré à 0 / à 1 \\
            & \V|supply0 supply1|& masse / alimentation \\
Variables   & \V|reg|            & affectée dans un process \\
            & \V|integer|        & entier 32 bits signé \\
            & \V|real|           & flottant \\
            & \V|time realtime|  & temps de simulation
\end{tab3}
\sub{Vecteurs et mémoires}
\begin{code}
wire [31:0] a;                  // vecteur 32 bits
reg  [15:0] b;                  // b[15] = MSB
integer t [0:100];              // 101 entiers
reg  [31:0] m [0:1023];         // 1024 mots de 32 bits
reg  [31:0] m2 [0:1023][1:20];  // tableau 2D
\end{code}
\sub{Conversions (cast) \sv}
\begin{code}
int'(3.7)        // = 4 (réel -> entier, arrondi)
8'(x)            // taille : tronque / étend
signed'(x)       // relit comme signé     ($signed)
unsigned'(x)     // relit comme non signé ($unsigned)
\end{code}
\end{bloc}
\end{colonne}\colbreak
% ======================= colonne 2 : types SystemVerilog =======================
\begin{colonne}[types SystemVerilog][col-b]
\begin{bloc}{Types SystemVerilog}{typ}
\begin{typetab}
\V|bit|      & 2 & variable & non signé & vecteur 2 états \\
\V|logic|    & 4 & variable & non signé & remplace \V|reg| / \V|wire| \\
\V|byte|     & 2 & 8        & signé     & ou caractère ASCII \\
\V|shortint| & 2 & 16       & signé     & \\
\V|int|      & 2 & 32       & signé     & \\
\V|longint|  & 2 & 64       & signé     & \\
\V|integer|  & 4 & 32       & signé     & \\
\V|time|     & 4 & 64       & non signé &
\end{typetab}
\sub{Chaînes : string}
\begin{grille}{3}
\V|s.len()|       & \V|s.getc(i)|     & \V|s.putc(i, c)| \\
\V|s.substr(i, j)| & \V|s.compare(s2)| & \V|s.toupper()| \\
\V|s.itoa(i)|     & \V|s.hextoa(i)|   & \V|s.atoi()|
\end{grille}
\end{bloc}

\begin{bloc}{Énumérations et structures}{agg}
\begin{code}
typedef enum {RED, YELLOW, GREEN} color_t;
typedef enum logic [1:0] {IDLE, RUN} state_t;  // codé
color_t c = RED;
c = c.next();  c.name();         // YELLOW, "YELLOW"
typedef struct packed {          // packed : un vecteur
  bit [7:0]  opcode;
  bit [23:0] addr;
} instr_t;
instr_t ir;
ir.opcode = 8'hAE;  ir = 32'hAE000010;   // champ · tout
typedef union packed { int i; bit [31:0] b; } mot_t;
\end{code}
\end{bloc}

\begin{bloc}{Tableaux}{agg}
\begin{code}
bit   [7:0] c1;              // packed : vecteur de bits
real  u [7:0];               // unpacked : tableau
int   t [8][32];             // = [0:7][0:31], notation C
logic [7:0] mem [0:255];     // mémoire 256 × 8 bits
data = mem[addr];            // lit le mot n°addr
A = B;  A[i:j] = B[i:j];  A == B;   // copie, tranche, ==
\end{code}
\sub{Dynamiques}
\begin{code}
bit [3:0] nib[];             // taille fixée au run
nib = new[12];               // alloue 12 éléments
nib = new[13](nib);          // agrandit, garde les 12
nib.size();  nib.delete();   // taille · vide tout
\end{code}
\sub{Files (queues)}
\begin{code}
int q[$] = {3, 2, 7};          // file initialisée
bit b[$:255];                  // 256 éléments au plus
q.push_back(1);  q.push_front(0);
x = q.pop_front();  y = q.pop_back();
q.insert(i, v);  q.delete(i);  q.size();
x = q[0];  y = q[$];           // premier · dernier
\end{code}
\sub{Tableaux associatifs}
\begin{code}
bit [20:0] tab[string];        // indexé par chaîne
tab["toto"] = 21'h13cafe;
if (tab.exists("titi")) tab.delete("titi");
foreach (tab[k]) $display("%s", k);  // clés triées
\end{code}
\note{Aussi : \V|.num()| \V|.first(k)| \V|.next(k)| \V|.last(k)| \V|.prev(k)| ; \V|int t[*]| : index quelconque.}
\end{bloc}
\end{colonne}\colbreak
% ======================= colonne 3 : contrôle =======================
\begin{colonne}[contrôle · temps][col-c]
\begin{bloc}{Structures de contrôle}{ctl}
\begin{code}
if («cond») begin
  «instructions»
end else if («cond2») begin
  «instructions»
end else begin
  «instructions»
end

case («expression»)
  «v1»      : «instr»;
  «v2»      : begin «instrs» end
  «v3», «v4»  : «instr»;      // même branche
  default : «instr»;      // évite les latchs
endcase

for (int i = 0; i < N; i++) begin ... end
repeat («n»)       begin ... end
while («cond»)     begin ... end
foreach (tab[i]) begin ... end
do begin ... end while («cond»);
forever begin ... end         // boucle infinie
\end{code}
\note{\V|casez| : \V|z| ou \V|?| = joker · \V|casex| : \V|x| et \V|z| jokers.}
\end{bloc}

\begin{bloc}{Délais et évènements}{ctl}
\begin{code}
`timescale 1ns/1ps        // unité / précision
#20;                      // attendre 20 unités
#10 a = 1;                // retarde l'instruction
wait («cond») #(«d») «instr»;   // attend cond vraie
@(«évènement») «instr»;       // attend l'évènement
@(posedge clk);           // front montant
@(negedge rst_n);         // front descendant
@(a or b)  @(a, b)  @*    // liste de sensibilité
event done;               // évènement nommé
-> done;                  // le déclenche
wait (done.triggered);    // l'attend (sans course)
\end{code}
\end{bloc}

\begin{bloc}{Tâches et fonctions}{ctl}
\begin{code}
task do_something;           // style Verilog
  input  REQ;
  output GNT;
  begin «instructions» end     // peut consommer du temps
endtask

task do_something(input  int          DATA,   // style SV
                  output logic [12:0] RESULT);
  «instructions»
endtask

function [15:0] max_value;   // style Verilog
  input [15:0] a, b;
  begin max_value = a < b ? b : a; end
endfunction

function bit [15:0] max_value(bit [15:0] a, b);
  return a < b ? b : a;      // style SV
endfunction
\end{code}
\note{\textbf{task} : peut contenir \V|#|, \V|@|, \V|wait| ; pas de valeur de retour.
\textbf{function} : s'exécute en temps nul, renvoie une valeur.}
\end{bloc}
\end{colonne}

%%%%%%%%%%%%%%%%%%%%%%%%%%%%%%%%%%%%%%%%%%%%%%%%%%%%%%%%%%%%%%%%%%%%%%%%%%%%%%
%                     VERSO : CONCEPTION ET SIMULATION
%%%%%%%%%%%%%%%%%%%%%%%%%%%%%%%%%%%%%%%%%%%%%%%%%%%%%%%%%%%%%%%%%%%%%%%%%%%%%%
\newpage
\entete{Verilog / SystemVerilog}{verso · conception et simulation ($\to$ vérif : fiche UVM)}
\noindent
% ======================= colonne 4 : module =======================
\begin{colonne}[module][col-a]
\begin{bloc}{Structure d'un module}{hw}
\begin{code}
module my_module_name (...);
  parameter «nom» = «valeur_défaut»;

  input  «signal»;
  output «signal»;

  assign «nom_wire» = «expression»;

  initial begin «instrs» end  // 0, 1 ou plusieurs
  always  begin «instrs» end  // 0, 1 ou plusieurs
  «task, function»

  generate
    genvar i;
    for (i = 0; i < 10; i++) begin : my_name
      «always, initial, instances, signaux»
    end
  endgenerate
endmodule
\end{code}
\sub{Instanciation}
\begin{code}
mod #(.N(8)) u1 (.clk(clk), .q(q)); // par nom (mieux)
mod u2 (clk, q);                    // par position
\end{code}
\end{bloc}

\begin{bloc}{Bloquant vs non bloquant}{hw}
\begin{optab}[6mm]
\V|=|  & \textbf{bloquant} : effectué avant l'instruction suivante \\
\V|<=| & \textbf{non bloquant} : mis à jour à la fin du pas de temps (delta suivant)
\end{optab}
\begin{code}
// Non bloquant : registre (séquentiel)
always @(posedge clk) begin
  a <= b;
  b <= a;                // échange a et b
end

// Bloquant en séquentiel : à éviter
always @(posedge clk) begin
  tmp = b;  b = a;  a = tmp;
end
\end{code}
\piege{séquentiel $\Rightarrow$ \V|<=|, combinatoire $\Rightarrow$ \V|=|. Jamais les deux dans un même bloc.}
\end{bloc}

\begin{bloc}{Mémoire (RAM, ROM)}{hw}
\begin{code}
logic [7:0] mem [0:255];
initial $readmemh("init.hex", mem);   // contenu initial
always_ff @(posedge clk) begin
  if (we) mem[addr] <= din;           // écriture
  dout <= mem[addr];                  // lecture synchrone
end
\end{code}
\note{\V|init.hex| : un mot hexa par ligne (\V|$readmemb| : binaire). Lecture synchrone et
écriture sur un front : reconnu comme bloc RAM à la synthèse.}
\end{bloc}

\begin{bloc}{Règles de synthèse}{hw}
\begin{optab}
\V|always_ff|   & séquentiel : \V|<=|, une horloge (+ reset asynchrone dans la liste) \\
\V|always_comb| & combinatoire : \V|=|, sortie affectée dans toutes les branches \\
\V|#| \V|initial| & ignorés (simulation seule ; \V|initial| admis sur FPGA) \\
Un pilote       & un signal n'est affecté que dans un seul \V|always| ou \V|assign|
\end{optab}
\end{bloc}
\end{colonne}\colbreak
% ======================= colonne 5 : circuits =======================
\begin{colonne}[circuits][col-b]
\begin{bloc}{Circuit combinatoire}{hw}
\begin{code}
module mux2 (input  logic a, b, sel,
             output logic f);
  always_comb
    if (sel == 0) f = a;
    else          f = b;
endmodule
\end{code}
\note{Équivalent : \V|assign f = sel ? b : a;|}
\piege{affecter la sortie dans \emph{toutes} les branches, sinon latch.}
\end{bloc}

\begin{bloc}{Élément séquentiel}{hw}
\begin{code}
module flipflop_syncr (d, q, clk, rst_n);
  input d, clk, rst_n;
  output logic q;
  always_ff @(posedge clk)
    if (rst_n == 0) q <= 0;  // reset synchrone
    else            q <= d;
endmodule
\end{code}
\note{Reset asynchrone : \V|always_ff @(posedge clk or negedge rst_n)|.}
\sub{Compteur avec enable}
\begin{code}
always_ff @(posedge clk)
  if (!rst_n)  cnt <= '0;       // '0 : tous les bits à 0
  else if (en) cnt <= cnt + 1;
\end{code}
\end{bloc}

\begin{bloc}{Machine à états (FSM)}{hw}
\begin{schema}[>={Stealth[length=1.6mm]}, node distance=14mm, line width=0.6pt,
    st/.style={circle, draw=trait, fill=paper, minimum size=7.5mm, inner sep=0pt,
               font=\ttfamily\bfseries\fontsize{6}{6}\selectfont},
    lbl/.style={font=\ttfamily\fontsize{6}{6}\selectfont, inner sep=1.5pt}]
  \node[st] (I) {IDLE};
  \node[st, right=of I, draw=cat, line width=1pt] (R) {RUN};
  \node[st, right=of R] (D) {DONE};
  \node[lbl, below=0.5mm of R, text=cat!80!black] {busy = 1};
  \draw[->] (I) -- node[lbl, above] {start} (R);
  \draw[->] (R) -- node[lbl, above] {stop} (D);
  \draw[->] (D.north) to[bend right=14] (I.north);
  \draw[->] (I) to[loop left, looseness=5] node[lbl, left] {!start} (I);
\end{schema}
\begin{code}
module fsm (input  logic clk, rst_n, start, stop,
            output logic busy);
  typedef enum logic [1:0] {IDLE, RUN, DONE} state_t;
  state_t state, next;

  // 1. registre d'état (séquentiel)
  always_ff @(posedge clk or negedge rst_n)
    if (!rst_n) state <= IDLE;
    else        state <= next;
  // 2. état suivant (combinatoire)
  always_comb begin
    next = state;                 // défaut : on reste
    case (state)
      IDLE    : if (start) next = RUN;
      RUN     : if (stop)  next = DONE;
      DONE    :            next = IDLE;
      default :            next = IDLE;
    endcase
  end
  // 3. sorties (Moore : fonction de l'état seul)
  assign busy = (state == RUN);
endmodule
\end{code}
\note{\textbf{Mealy} : sorties fonction de l'état \emph{et} des entrées $\to$ les calculer dans le \V|always_comb|.
\V|next = state| par défaut évite les latchs.}
\end{bloc}
\end{colonne}\colbreak
% ======================= colonne 6 : simulation =======================
\begin{colonne}[simulation][col-c]
\begin{bloc}{Testbench}{tb}
\begin{code}
`timescale 1ns/100ps
module test_adder;
  logic [3:0] a, b;                  // signaux internes
  logic [4:0] s;
  adder dut (.a(a), .b(b), .s(s));   // instanciation

  logic clk = 1'b1;
  always #5 clk = ~clk;              // période 10 ns

  initial begin
    a = 0;  b = 3;                   // à 0 ns
    #10 a = 1;                       // à 10 ns
    @(posedge clk) a = 2;            // au front suivant
    $display("t=%0t s=%h", $time, s);
    $finish;                         // fin de simulation
  end
endmodule
\end{code}
\note{\V|$monitor| affiche à chaque changement · \V|$stop| pause · \V|$urandom| aléatoire.}
\sub{Formats de \$display}
\begin{tab3}[0mm]{30mm}
 & \V|%b %o %d %h| & binaire, octal, décimal, hexa \\
 & \V|%s %t %0d|   & chaîne, temps, sans espaces
\end{tab3}
\end{bloc}

\begin{bloc}{Interfaces}{tb}
\begin{code}
interface apb_if (input logic clk);
  logic        psel, penable, pready;
  logic [31:0] paddr;
  modport slave  (input  psel, penable, paddr,
                  output pready);   // master : inverse
  clocking cb @(posedge clk);     // vue synchrone
    default input #1step output #1;
    output psel, paddr;           // front + #1
    input  pready;                // juste avant
  endclocking
endinterface
module dut (apb_if.slave s); «...s.psel...» endmodule
apb_if bus (clk);                 // instance (top)
dut u (.s(bus));
virtual apb_if vif = bus;         // handle (classes)
@(vif.cb);  vif.cb.psel <= 1;     // synchrone
\end{code}
\end{bloc}

\begin{bloc}{Processus parallèles}{tb}
\begin{code}
fork
  «instruction»;            // branches lancées
  begin «a»; «b»; end       // en parallèle
join                      // attend toutes
// join_any : la 1re · join_none : aucune
wait fork;                // attend les fils restants
disable fork;             // tue les fils restants
\end{code}
\sub{Communication entre processus}
\begin{code}
mailbox #(string) mb = new();  // FIFO, new(N) : bornée
mb.put("hi");  mb.get(s);      // bloquants
mb.try_get(s); mb.num();       // non bloquant · taille
semaphore sem = new(1);        // 1 clé
sem.get(1); «section critique» sem.put(1);
\end{code}
\end{bloc}
\end{colonne}
\end{document}
